% unit U3

% unit macros
\ifcsname remark\endcsname\else
  \theoremstyle{acmdefinition}\newtheorem{remark}[theorem]{Remark}
  \theoremstyle{acmplain}
\fi

\FloatBarrier
\subsection{Stores and the machine}\label{sec:mini-machine}

A program runs against a store. The store holds the capture binders of the platform, the values the program starts with, and the values the run allocates.

\begin{figure}[t]
\small
\color{tgtcol}
\begin{mathpar}
\inferrule[S-Nil]{ }{\storej{\cdot}{\cdot}}
\and
\inferrule[S-Val]{\storej{\sigma}{\Gamma} \\ \valj{\Gamma}{v}{T}}{\storej{\sigma, x = v}{\Gamma, x : T}}
\and
\inferrule[S-Cap]{\storej{\sigma}{\Gamma}}{\storej{\sigma, \kappa}{\Gamma, \kappa : \cstar}}
\\
\inferrule[K-Nil]{ }{\contj{\Gamma}{\cdot}{T}{T}}
\and
\inferrule[K-Let]{\tmj{\Gamma, x : T}{u}{\weak{T'}} \\ \clej{\Gamma, x : T}{f}{\uses{u}}{\weak{U}} \\ \contj{\Gamma}{K}{T'}{T''}}{\contj{\Gamma}{K \triangleright \frlet{x}{u}{U}{f}}{T}{T''}}
\\
\inferrule[St]{\storej{\sigma}{\Gamma} \\ \tmj{\Gamma}{t}{T} \\ \contj{\Gamma}{K}{T}{T'}}{\statej{\mst{\sigma}{K}{t}}{T'}}
\end{mathpar}
\[
\begin{array}{@{}p{5.5em}@{}l@{\enspace}c@{\enspace}l@{}}
\rn{E-Let} & \mst{\sigma}{K}{\letc{x}{t}{u}{U}{f}} & \step & \mst{\sigma}{K \triangleright \frlet{x}{u}{U}{f}}{t} \\
\rn{E-Alloc} & \mst{\sigma}{K \triangleright \frlet{x}{u}{U}{f}}{v} & \step & \mst{(\sigma, x = v)}{K}{u} \\
& & & \quad \text{if } x \text{ is not bound in } \sigma \\
\rn{E-Rename} & \mst{\sigma}{K \triangleright \frlet{x}{u}{U}{f}}{y} & \step & \mst{\sigma}{K}{u[x \mapsto y]} \\
\rn{E-App} & \mst{\sigma}{K}{x\,y} & \step & \mst{\sigma}{K}{t[z \mapsto y]} \\
& & & \quad \text{if } \sigma(x) = \lamc{A}{z}{T}{t}{g}
\end{array}
\]
\begin{mathpar}
\inferrule[Run-Refl]{ }{\mathit{st} \steps \mathit{st}}
\and
\inferrule[Run-Step]{\mathit{st} \steps \mathit{st}' \\ \mathit{st}' \step \mathit{st}''}{\mathit{st} \steps \mathit{st}''}
\end{mathpar}
\[
\begin{array}{@{}p{5.5em}@{}l@{}}
\text{Final state} & \mst{\sigma}{\cdot}{v} \\
\text{Final state} & \mst{\sigma}{\cdot}{x}
\end{array}
\]
\color{black}
\caption{Stores, continuations, states and the machine of the fragment. $\contj{\Gamma}{K}{T}{T'}$ says that $K$ accepts a value of type $T$ and produces one of type $T'$.}
\label{fig:mini-machine}
\end{figure}

\paragraph{Stores.}
Figure~\ref{fig:mini-machine} gives the rules. \rn{S-Cap} binds a capture variable. \rn{S-Val} types a value in the context of the entries before it and binds its location at the type of the value. So the store is inert in the sense of report~1. The type of a location is the type of its value and asserts nothing more. In particular, by \rn{V-Lam} and \rn{V-Unit}, the capture set of that type is the annotation of the value. The rules are syntax directed, and a value has one type in a given context. So a typed store has exactly one context, and we call it the context of the store.

\paragraph{The machine.}
A state $\mst{\sigma}{K}{t}$ has a store, a continuation of let frames, and a term. \rn{E-Let} pushes the body of a let as a frame and runs the bound term. When the bound term has become a value, \rn{E-Alloc} stores it at the variable of the frame and runs the body. When it has become a variable $y$, \rn{E-Rename} puts $y$ for the variable of the frame in the body, in its capture sets and evidence too. \rn{E-App} reads the closure stored at the function of a call and runs its body at the argument. These are the steps of report~1's machine for lets and for application through a variable (Figure~\ref{R1-fig:mini-machine} of report~1), without casts. No step reads the brackets. Each step keeps the store or appends one location.

\paragraph{Continuations and states.}
\rn{K-Let} has the last two premises of \rn{T-Let}, so a frame carries the avoidance evidence of the let it came from. \rn{St} types a state when its term and its continuation are typed in the context of its store. The evidence of a frame is typed in that context extended by the variable of the frame. A state is typed without a use set. Its use set is that of its term and the sets its frames declare (Figure~\ref{fig:mini-typing}), as Idea~1 promised.

\paragraph{The running example.}
By \rn{S-Cap} and \rn{S-Val}, $\storej{\sigma_0}{\Gamma_1}$, and \rn{K-Nil} and \rn{St} type both initial states at $\Unit$. Table~\ref{tab:mini-run} runs $\mathit{prog}_1$, with
\[
\begin{array}{@{}lcl@{\qquad}lcl@{}}
K_0 & = & \cdot \triangleright \frlet{\mathit{inner}}{\mathit{inner}\,\mathit{unit}}{\cset{\mathit{console}, \mathit{log}}}{f_1} &
\sigma_1 & = & \sigma_0,\ \mathit{inner} = v_{\mathit{inner}} \\
K_1 & = & K_0 \triangleright \frlet{r}{v_{\mathit{inner}}}{\cempty}{\ev{elem}\,\cempty\,\cempty} &
T_{\mathit{inner}} & = & \capt{(\pit[z]{\Unit}{\Unit})}{\cset{\mathit{console}}} \\
t_2 & = & \letc{r}{\mathit{log}\,\mathit{unit}}{v_{\mathit{inner}}}{\cempty}{\ev{elem}\,\cempty\,\cempty} &
\Gamma_2 & = & \Gamma_1,\ \mathit{inner} : T_{\mathit{inner}}
\end{array}
\]
\begin{table}[t]
\footnotesize
\setlength{\tabcolsep}{2.5pt}
\begin{tabular}{@{}c l c c l l l l p{24mm}@{}}
\toprule
$k$ & term & $K$ & $\sigma$ & use set & roots & inspects & step & \raggedright case of Lemma~\ref{lem:mini-step} \tabularnewline
\midrule
0 & $\mathit{prog}_1$ & $\cdot$ & $\sigma_0$ & $\cset{\mathit{c1}, \mathit{unit}, \mathit{console}, \mathit{log}}$ & $\cset{\kappa_1, \kappa_2}$ & none & \rn{E-Let} & \raggedright equal sets \tabularnewline
1 & $\mathit{c1}\,\mathit{unit}$ & $K_0$ & $\sigma_0$ & $\cset{\mathit{c1}, \mathit{unit}, \mathit{console}, \mathit{log}}$ & $\cset{\kappa_1, \kappa_2}$ & $\mathit{c1}$ ($\kappa_1$) & \rn{E-App} & \raggedright closing evidence of $v_{\mathit{c1}}$ at $\mathit{unit}$ \tabularnewline
2 & $t_2$ & $K_0$ & $\sigma_0$ & $\cset{\mathit{log}, \mathit{unit}, \mathit{console}}$ & $\cset{\kappa_1, \kappa_2}$ & none & \rn{E-Let} & \raggedright equal sets \tabularnewline
3 & $\mathit{log}\,\mathit{unit}$ & $K_1$ & $\sigma_0$ & $\cset{\mathit{log}, \mathit{unit}, \mathit{console}}$ & $\cset{\kappa_1, \kappa_2}$ & $\mathit{log}$ ($\kappa_1$) & \rn{E-App} & \raggedright closing evidence of $v_{\mathit{log}}$ at $\mathit{unit}$ \tabularnewline
4 & $\mathit{unit}$ & $K_1$ & $\sigma_0$ & $\cset{\mathit{unit}, \mathit{console}, \mathit{log}}$ & $\cset{\kappa_1, \kappa_2}$ & none & \rn{E-Rename} & \raggedright avoidance evidence $\ev{elem}\,\cempty\,\cempty$ at $\mathit{unit}$ \tabularnewline
5 & $v_{\mathit{inner}}$ & $K_0$ & $\sigma_0$ & $\cset{\mathit{console}, \mathit{log}}$ & $\cset{\kappa_1, \kappa_2}$ & none & \rn{E-Alloc} & \raggedright avoidance evidence $f_1$ over $\sigma_1$ \tabularnewline
6 & $\mathit{inner}\,\mathit{unit}$ & $\cdot$ & $\sigma_1$ & $\cset{\mathit{inner}, \mathit{unit}}$ & $\cset{\kappa_2}$ & $\mathit{inner}$ ($\kappa_2$) & \rn{E-App} & \raggedright closing evidence of $v_{\mathit{inner}}$ at $\mathit{unit}$ \tabularnewline
7 & $\mathit{console}\,\mathit{unit}$ & $\cdot$ & $\sigma_1$ & $\cset{\mathit{console}, \mathit{unit}}$ & $\cset{\kappa_2}$ & $\mathit{console}$ ($\kappa_2$) & & \tabularnewline
\bottomrule
\end{tabular}
\caption{The run of $\mathit{prog}_1$. Row $k$ is state $k$ and the step taken from it. Use sets are shown as sets, without duplicates. An inspected variable is shown with its roots. Roots are computed in $\Gamma_1$ for rows~0 to~5 and in $\Gamma_2$ for rows~6 and~7. The last column names the case of Lemma~\ref{lem:mini-step} that bounds the next row by this one. The table stops at state~7, which inspects $\mathit{console}$.}
\label{tab:mini-run}
\end{table}
\rn{E-Let} pushes $K_0$, and the state calls $\mathit{c1}$ (row~1). \rn{E-App} runs the body of $v_{\mathit{c1}}$ at $\mathit{unit}$, which is $t_2$. Its let pushes the frame of $K_1$ and calls $\mathit{log}$ (row~3), which returns its argument. \rn{E-Rename} puts $\mathit{unit}$ for $r$ in $v_{\mathit{inner}}$, which does not mention $r$ (row~5). \rn{E-Alloc} stores $v_{\mathit{inner}}$ at $\mathit{inner}$. The context of the new store $\sigma_1$ is $\Gamma_2$, since $v_{\mathit{inner}}$ has type $T_{\mathit{inner}}$ in $\Gamma_1$. The body $\mathit{inner}\,\mathit{unit}$ calls the stored closure, which calls $\mathit{console}$ (row~7). One more step by \rn{E-App} returns $\mathit{unit}$ in a final state. Status: the typings of $\sigma_0$, of both initial states and of $\sigma_1$ in $\Gamma_2$ are decided. The seven steps of the table and the root $\kappa_2$ of $\mathit{console}$ in row~7 are built by hand. The eighth step and the other roots are on paper.

\FloatBarrier
\subsection{One step}\label{sec:mini-step}

Preservation has the form it has in report~1. What is new is its use-set half: one step never adds to the roots of the use set of a state. Both halves read the typed store through one lemma and move evidence by substitution, which comes first.

\begin{lemma}[Substitution and weakening]\label{lem:mini-subst}
\leavevmode
\begin{enumerate}
\item[(a)] If $\clej{\Gamma, x : T}{f}{C}{D}$ and $\Gamma(y) = T$, then $\clej{\Gamma}{f[x \mapsto y]}{C[x \mapsto y]}{D[x \mapsto y]}$.
\item[(b)] If $\tmj{\Gamma, x : T}{t}{T'}$ and $\Gamma(y) = T$, then $\tmj{\Gamma}{t[x \mapsto y]}{T'[x \mapsto y]}$. The same holds for the typing of values.
\item[(c)] $\uses{t[x \mapsto y]} = \uses{t}[x \mapsto y]$.
\item[(d)] A typing of capture evidence, of a term, of a value or of a continuation in $\Gamma$ holds in $\Gamma, \Gamma'$ for any binders $\Gamma'$ whose variables are fresh.
\end{enumerate}
\end{lemma}

\begin{proof}[Proof sketch]
Parts (a), (b) and (d) go by induction on the derivation, and (c) by induction on $t$. Since \rn{T-Let}, \rn{V-Lam} and \rn{K-Let} extend the context, the induction proves (a) and (b) for $\Gamma, x : T, \Delta$ with the result in $\Gamma, \Delta[x \mapsto y]$, and (d) for fresh binders inserted before $\Delta$. The one case of (a) that looks at $x$ is \rn{C-Var} at $x$ itself. There $\ev{capvar}\,x$ proves $\cset{x} \le C_T$ for the capture set $C_T$ of $T$. Since $\Gamma(y) = T$ and $C_T$ does not mention $x$, $\ev{capvar}\,y$ proves $\cset{y} \le C_T$.
\end{proof}

\begin{lemma}[What a typed store says]\label{lem:mini-store}
Let $\storej{\sigma}{\Gamma}$ and $\sigma(x) = v$.
\begin{enumerate}
\item[(a)] $\Gamma(x) = \capt{S}{\annot{v}}$ for some $S$. Hence $\capeqv{\Gamma}{\cset{x}}{\annot{v}}$.
\item[(b)] If $v = \lamc{A}{z}{T}{t}{g}$, then $\Gamma(x) = \capt{(\pit[z]{T}{T'})}{A}$ for some $T'$, $\tmj{\Gamma, z : T}{t}{T'}$ and $\clej{\Gamma, z : T}{g}{\uses{t}}{\weak{A} \cup \cset{z}}$.
\end{enumerate}
\end{lemma}

\begin{proof}
Let $\sigma = \sigma', x = v, \sigma''$. Inversion of the store typing gives $\Gamma = \Gamma', x : T_0, \Gamma''$ with $\valj{\Gamma'}{v}{T_0}$, and weakening, Lemma~\ref{lem:mini-subst}(d), gives $\valj{\Gamma}{v}{T_0}$. Inversion of \rn{V-Lam} or \rn{V-Unit} in $\Gamma$ shows that the capture set of $T_0$ is $\annot{v}$, and for a closure it gives the premises of (b). Lemma~\ref{lem:mini-var} gives (a).
\end{proof}

\begin{theorem}[Preservation]\label{thm:mini-pres}
If $\statej{\mathit{st}}{T}$ and $\mathit{st} \step \mathit{st}'$, then $\statej{\mathit{st}'}{T}$.
\end{theorem}

\begin{proof}[Proof sketch]
As in Section~\ref{R1-sec:mini-safety} of report~1, by cases on the step, moving evidence by Lemma~\ref{lem:mini-subst}. \rn{E-Alloc} stores the value at the type its frame expects, so \rn{S-Val} extends the context of the store by exactly the binder under which \rn{K-Let} typed the body. \rn{E-App} takes the body of the closure from Lemma~\ref{lem:mini-store}(b), and \rn{T-App} gives the argument the parameter type.
\end{proof}

\begin{lemma}[The use-set half of preservation]\label{lem:mini-step}
Let $\mathit{st}$ be a state with store $\sigma$, let $\storej{\sigma}{\Gamma}$ and $\statej{\mathit{st}}{T}$, and let $\mathit{st} \step \mathit{st}'$, where the store $\sigma'$ of $\mathit{st}'$ has $\storej{\sigma'}{\Gamma'}$. Then $\sext{\sigma}{\sigma'}$ and $\caple{\Gamma'}{\uses{\mathit{st}'}}{\uses{\mathit{st}}}$.
\end{lemma}

\begin{proof}[Proof sketch]
Every step keeps the store or appends one location, so $\sext{\sigma}{\sigma'}$. When it keeps the store, $\Gamma' = \Gamma$, since a store has one context. By \rn{St}, the term and the continuation of $\mathit{st}$ are typed in $\Gamma$.
\begin{itemize}
\item \rn{E-Let} (rows~0 and~2). Before: $\uses{t} \cup U \cup \uses{K}$. After: $\uses{t} \cup \uses{K} \cup U$. The two sets have the same atoms.
\item \rn{E-Rename} (row~4). Before: $\cset{y} \cup \uses{K} \cup U$. After: $\uses{u[x \mapsto y]} \cup \uses{K}$. \rn{K-Let} gives $\clej{\Gamma, x : T_0}{f}{\uses{u}}{\weak{U}}$ for the input type $T_0$ of the frame, and \rn{T-Var} gives $\Gamma(y) = T_0$. Lemma~\ref{lem:mini-subst}(a) and~(c) give $\clej{\Gamma}{f[x \mapsto y]}{\uses{u[x \mapsto y]}}{U}$, since $U$ does not mention $x$. Theorem~\ref{thm:mini-cap} gives $\caple{\Gamma}{\uses{u[x \mapsto y]}}{U}$.
\item \rn{E-Alloc} (row~5). Before: $\uses{K} \cup U$, since a value uses nothing. After: $\uses{u} \cup \uses{K}$. The value has the input type $T_0$ of the frame, so the context of $\sigma, x = v$ is $\Gamma' = \Gamma, x : T_0$. That is the context in which \rn{K-Let} typed the avoidance evidence $f$ of the frame, for $\uses{u} \le \weak{U}$. So $f$ is now closed evidence, and Theorem~\ref{thm:mini-cap} in $\Gamma'$ gives $\caple{\Gamma'}{\uses{u}}{U}$.
\item \rn{E-App} (rows~1, 3 and~6). Before: $\cset{x, y} \cup \uses{K}$. After: $\uses{t[z \mapsto y]} \cup \uses{K}$, where $\sigma(x) = \lamc{A}{z}{T}{t}{g}$. Lemma~\ref{lem:mini-store}(b) gives $\Gamma(x) = \capt{(\pit[z]{T}{T'})}{A}$ and $\clej{\Gamma, z : T}{g}{\uses{t}}{\weak{A} \cup \cset{z}}$, and \rn{T-App} gives $\Gamma(y) = T$. Lemma~\ref{lem:mini-subst}(a) and~(c) give $\clej{\Gamma}{g[z \mapsto y]}{\uses{t[z \mapsto y]}}{A \cup \cset{y}}$, and Theorem~\ref{thm:mini-cap} gives $\caple{\Gamma}{\uses{t[z \mapsto y]}}{A \cup \cset{y}}$. Lemma~\ref{lem:mini-store}(a) gives $\capeqv{\Gamma}{\cset{x}}{A}$, so $\caple{\Gamma}{A \cup \cset{y}}{\cset{x, y}}$. \qedhere
\end{itemize}
\end{proof}

Each case reads evidence that a frame or a stored closure carries, and no derivation of capture evidence is inverted. The last column of Table~\ref{tab:mini-run} names these cases. Each use of Theorem~\ref{thm:mini-cap} is at closed evidence, typed in the context of the current store. That is the form in which Section~\ref{sec:target-canon} proves the theorem once capture members enter.

\begin{example}[The allocation step]\label{ex:mini-alloc}
State~5 is $\mst{\sigma_0}{K_0}{v_{\mathit{inner}}}$. Its use set is $\cset{\mathit{console}, \mathit{log}}$, the set that $K_0$ declares, with roots $\cset{\kappa_1, \kappa_2}$ in $\Gamma_1$. \rn{E-Alloc} stores $v_{\mathit{inner}}$ at $\mathit{inner}$ and runs $\mathit{inner}\,\mathit{unit}$. The avoidance evidence $f_1$ of the frame was typed in $\Gamma_1, \mathit{inner} : T_{\mathit{inner}}$ when $\mathit{prog}_1$ was checked, with $\mathit{inner}$ the variable of a let (Example~\ref{ex:mini-typing}). That context is now $\Gamma_2$, the context of $\sigma_1$, so $f_1$ is closed evidence for $\cset{\mathit{inner}, \mathit{unit}} \le \cset{\mathit{console}, \mathit{log}}$. By Lemma~\ref{lem:mini-store}(a) the roots of $\mathit{inner}$ are those of the annotation $\cset{\mathit{console}}$ of the stored closure, $\cset{\kappa_2}$. Theorem~\ref{thm:mini-cap} puts them among $\cset{\kappa_1, \kappa_2}$. The roots of the state drop to $\cset{\kappa_2}$. The frame declared $\mathit{log}$, which its body never calls, and the step forgets it. Status: the step and the roots of $\cset{\mathit{console}}$ in $\Gamma_2$ are built by hand, the typing of $\sigma_1$ is decided, and the rest is on paper.
\end{example}

\FloatBarrier
\subsection{Along a run}\label{sec:mini-run}

Along a run the store grows, and roots are computed in a new context after each allocation. One lemma relates the old roots to the new ones.

\begin{lemma}[Roots under store extension]\label{lem:mini-ext}
Let $\storej{\sigma}{\Gamma}$, $\storej{\sigma'}{\Gamma'}$ and $\sext{\sigma}{\sigma'}$. Then $\Gamma' = \Gamma, \Gamma''$ for some $\Gamma''$, and $\roots{\Gamma'}{C} = \roots{\Gamma}{C}$ for every capture set $C$ whose atoms are bound in $\Gamma$. Hence, for such $C$ and $D$, $\caple{\Gamma}{C}{D}$ gives $\caple{\Gamma'}{C}{D}$, and $\kappa$ is a root of $C$ in $\Gamma'$ if and only if it is one in $\Gamma$.
\end{lemma}

\begin{proof}
Let $\sigma' = \sigma, \sigma''$. Store typing types each entry in the context of the entries before it, so the typing of $\sigma'$ types $\sigma$ as before, and $\Gamma' = \Gamma, \Gamma''$. The equations for roots read the binder of an atom and the part of the context before it. For an atom bound in $\Gamma$, both are the same in $\Gamma'$.
\end{proof}

\begin{theorem}[Capture prediction]\label{thm:mini-pred}
Let $\statej{\mathit{st}}{T}$ and $\mathit{st} \steps \mathit{st}'$, let $\sigma$ and $\sigma'$ be the stores of $\mathit{st}$ and $\mathit{st}'$, and let $\storej{\sigma'}{\Gamma'}$. Then $\sext{\sigma}{\sigma'}$ and $\caple{\Gamma'}{\uses{\mathit{st}'}}{\uses{\mathit{st}}}$.
\end{theorem}

\begin{proof}
By induction on the run. For \rn{Run-Refl} both claims hold by reflexivity. For \rn{Run-Step}, let $\mathit{st} \steps \mathit{st}_m \step \mathit{st}'$, and let $\sigma_m$ be the store of $\mathit{st}_m$. Theorem~\ref{thm:mini-pres} along the run gives $\statej{\mathit{st}_m}{T}$, so $\sigma_m$ has a context $\Gamma_m$. The induction hypothesis gives $\sext{\sigma}{\sigma_m}$ and $\caple{\Gamma_m}{\uses{\mathit{st}_m}}{\uses{\mathit{st}}}$. Lemma~\ref{lem:mini-step} gives $\sext{\sigma_m}{\sigma'}$ and $\caple{\Gamma'}{\uses{\mathit{st}'}}{\uses{\mathit{st}_m}}$. The atoms of both use sets are bound in $\Gamma_m$, since both states are typed and the context of $\sigma$ is a prefix of $\Gamma_m$ (Lemma~\ref{lem:mini-ext}). So Lemma~\ref{lem:mini-ext} carries the first inclusion into $\Gamma'$, and the two inclusions compose. So do the two extensions.
\end{proof}

\begin{remark}[An inspected variable is used]\label{rem:mini-insp}
If $\insp{\mathit{st}} = x$, then $x$ is an atom of $\uses{\mathit{st}}$, so $\caple{\Gamma}{\cset{x}}{\uses{\mathit{st}}}$ for every $\Gamma$. Only a call $x\,y$ inspects, and its use set contains $x$.
\end{remark}

\begin{theorem}[Effect safety]\label{thm:mini-effect}
Let $\statej{\mathit{st}}{T}$, let $\Gamma$ be the context of the store of $\mathit{st}$, and let $\kappa$ be a capture variable bound in $\Gamma$ that is not a root of $\uses{\mathit{st}}$ in $\Gamma$. Let $\mathit{st} \steps \mathit{st}'$, where the store $\sigma'$ of $\mathit{st}'$ has $\storej{\sigma'}{\Gamma'}$, and let $\insp{\mathit{st}'} = x$. Then $\kappa \notin \roots{\Gamma'}{\cset{x}}$.
\end{theorem}

\begin{proof}
Suppose $\kappa \in \roots{\Gamma'}{\cset{x}}$. Remark~\ref{rem:mini-insp} makes $\kappa$ a root of $\uses{\mathit{st}'}$ in $\Gamma'$, Theorem~\ref{thm:mini-pred} a root of $\uses{\mathit{st}}$ in $\Gamma'$, and Lemma~\ref{lem:mini-ext} a root of $\uses{\mathit{st}}$ in $\Gamma$, since the atoms of $\uses{\mathit{st}}$ are bound in $\Gamma$. That contradicts the hypothesis on $\kappa$.
\end{proof}

For $\mathit{prog}_1$ the theorem says nothing about $\kappa_2$, a root of its initial use set. The run agrees with capture prediction. State~7 inspects $\mathit{console}$, whose root is $\kappa_2$, and Remark~\ref{rem:mini-insp} with Theorem~\ref{thm:mini-pred} puts $\kappa_2$ among the roots of the initial use set (built by hand).

\begin{example}[The second program]\label{ex:mini-prog2}
The state $\mst{\sigma_0}{\cdot}{\mathit{prog}_2}$ is typed, and $\kappa_2$ is not a root of $\uses{\mathit{prog}_2}$, whose roots in $\Gamma_1$ are $\cset{\kappa_1}$. By Theorem~\ref{thm:mini-effect}, no state reached from $\mathit{prog}_2$ inspects a variable with root $\kappa_2$. The Lean proves this instance for every such state (built by hand). On paper, the run applies the rules of the first six steps of Table~\ref{tab:mini-run} to the frames of $\mathit{prog}_2$. At the allocation the frame declares $\cempty$, and $\ev{capvar}\,\mathit{unit}$ proves $\cset{\mathit{unit}} \le \cempty$. So the closure $v_{\mathit{inner}}$ enters the store at $\mathit{inner}$ and never a use set. The run inspects $\mathit{c1}$ and $\mathit{log}$ and ends in the final state $\mst{\sigma_1}{\cdot}{\mathit{unit}}$. It never calls $\mathit{console}$, nor the stored closure $\mathit{inner}$ that would. An instance of effect safety is informative when its run inspects a variable $x$ for which $\cset{x}$ has a root. Otherwise its conclusion holds for every $\kappa$ and says nothing. This one is informative, since $\mathit{c1}$ and $\mathit{log}$ have the root $\kappa_1$ (on paper).
\end{example}

\FloatBarrier
\subsection{What the miniature leaves out}\label{sec:mini-out}

The full calculus adds the following to the fragment, and Sections~\ref{sec:target} to~\ref{sec:classifiers} explain each.
\begin{itemize}
\item Capture members, and the capture names $\sel{x}{\ell}$ that select them. They make bad capture bounds derivable, and they bring the store into the theorem about capture evidence.
\item Report~1's shape evidence and casts, with the application steps that read a chain of casts.
\item Boxes, which hide a capture.
\item Progress and the erasure simulation, which keep the form they have in report~1.
\item The source and its translation, scopes and levels, and classifiers.
\end{itemize}
None of these changes the shape of the argument: one theorem about capture evidence, and a use-set lemma whose cases read the evidence that the state carries.
